\documentclass[aspectratio=169,12pt]{beamer}

\usepackage{fontspec}
\setsansfont{DejaVu Sans}
\setmonofont{DejaVu Sans Mono}
\usepackage{graphicx}
\usepackage{booktabs}
\usepackage{tabularx}
\usepackage{amsmath}
\usepackage{tikz}
\usepackage{xcolor}
\usepackage{listings}
\usepackage{hyperref}

\usetheme{AnnArbor}
\usefonttheme{professionalfonts}
\setbeamertemplate{navigation symbols}{}

\definecolor{Navy}{HTML}{173B57}
\definecolor{Blue}{HTML}{1F5AA6}
\definecolor{Teal}{HTML}{2A9D8F}
\definecolor{Orange}{HTML}{F28E2B}
\definecolor{Red}{HTML}{D1495B}
\definecolor{Ink}{HTML}{17202A}
\definecolor{Muted}{HTML}{667085}
\definecolor{Paper}{HTML}{F7F9FC}
\definecolor{CodeBg}{HTML}{F1F4F8}

\setbeamercolor{alerted text}{fg=Red}
\setbeamerfont{frametitle}{series=\bfseries,size=\large}
\setbeamerfont{title}{series=\bfseries,size=\LARGE}
\setbeamerfont{subtitle}{size=\large}

\lstdefinestyle{py}{
  language=Python,
  basicstyle=\ttfamily\fontsize{9.3}{11.2}\selectfont,
  keywordstyle=\color{Blue}\bfseries,
  stringstyle=\color{Teal},
  commentstyle=\color{Muted}\itshape,
  numberstyle=\tiny\color{Muted},
  numbers=left,
  numbersep=5pt,
  backgroundcolor=\color{CodeBg},
  frame=single,
  rulecolor=\color{CodeBg},
  breaklines=true,
  showstringspaces=false,
  columns=fullflexible,
  keepspaces=true,
  aboveskip=2pt,
  belowskip=2pt
}

\newcommand{\asset}[1]{figures_annarbor/#1}
\newcommand{\key}[1]{\textcolor{Blue}{\textbf{#1}}}
\newcommand{\muted}[1]{\textcolor{Muted}{#1}}
\newcommand{\good}[1]{\textcolor{Teal}{\textbf{#1}}}
\newcommand{\bad}[1]{\textcolor{Red}{\textbf{#1}}}

\newenvironment{tightitems}{\begin{itemize}\setlength\itemsep{5pt}}{\end{itemize}}

\newcommand{\SectionSlide}[2]{
\section{#1}
\begin{frame}{#1}
  \vfill
  \centering
  {\usebeamercolor[fg]{structure}\fontsize{24}{29}\selectfont\bfseries #1\par}
  \vspace{5mm}
  {\large #2\par}
  \vfill
\end{frame}}

\title{Trực quan hóa dữ liệu với Matplotlib}
\subtitle{Từ biểu đồ cơ bản đến quy trình phân tích dữ liệu có thể tái lập}
\author{Duc-Minh Vu}
\institute{SLSCM Lab -- FDA -- NEU}
\date{\today}

\begin{document}

% 1
\begin{frame}
  \titlepage
\end{frame}

% 2
\begin{frame}{Một biểu đồ tốt giúp trả lời câu hỏi, không chỉ “làm đẹp” dữ liệu}
\begin{columns}[T,onlytextwidth]
\column{.49\textwidth}
\begin{block}{Bảng số liệu cho biết}
12 giá trị doanh thu theo tháng, chính xác đến từng con số.
\end{block}
\begin{block}{Biểu đồ cho thấy}
Xu hướng tăng, giai đoạn chững lại và mùa cao điểm cuối năm.
\end{block}
\column{.49\textwidth}
\includegraphics[width=\linewidth]{\asset{line_revenue.png}}
\end{columns}
\vspace{1mm}
\begin{center}\key{Thông điệp:} trực quan hóa nén dữ liệu thành mẫu hình mà con người nhận ra nhanh.\end{center}
\end{frame}

% 3
\begin{frame}{Mục tiêu học tập}
Sau bài học, người học có thể:
\begin{tightitems}
  \item giải thích mô hình \key{Figure--Axes} và chọn đúng API;
  \item tạo line, bar, scatter, histogram và pie chart;
  \item tùy biến tiêu đề, nhãn, ticks, legend, màu, marker và grid;
  \item bố trí nhiều biểu đồ bằng \texttt{subplots};
  \item tạo heatmap, contour, 3D và animation ở mức nhập môn;
  \item tích hợp Matplotlib với NumPy, Pandas và Seaborn;
  \item xuất hình chất lượng cao và tránh các lỗi trình bày gây hiểu sai.
\end{tightitems}
\end{frame}

% 4
\begin{frame}{Lộ trình bài học đi từ câu hỏi đến sản phẩm trực quan}
\begin{enumerate}\setlength\itemsep{8pt}
  \item Matplotlib là gì và nằm ở đâu trong hệ sinh thái Python?
  \item Figure, Axes và hai phong cách viết mã.
  \item Năm loại biểu đồ nền tảng.
  \item Tùy biến, chú thích và bố cục nhiều biểu đồ.
  \item Heatmap, 3D, contour, animation và tương tác.
  \item Pandas/Seaborn, lưu hình và quy trình Data Science.
  \item Quiz, bài tập có hướng dẫn và bài tập tự làm.
\end{enumerate}
\end{frame}

\SectionSlide{1. Nền tảng Matplotlib}{Hiểu đúng thư viện trước khi viết mã}

% 6
\begin{frame}{Matplotlib là lớp nền trực quan hóa của hệ sinh thái Python}
\begin{columns}[T,onlytextwidth]
\column{.56\textwidth}
\begin{tightitems}
  \item Thư viện mã nguồn mở cho hình \key{tĩnh}, \key{động} và \key{tương tác}.
  \item Làm việc tự nhiên với \texttt{NumPy ndarray} và \texttt{Pandas DataFrame}.
  \item Kiểm soát chi tiết trục, nhãn, legend, màu, font và layout.
  \item Xuất PNG, PDF, SVG và nhiều định dạng khác.
  \item Có thể nhúng vào Tkinter, Qt, GTK hoặc wxPython.
\end{tightitems}
\column{.40\textwidth}
\begin{block}{Vai trò thực tế}
\textbf{Khám phá:} tìm lỗi, ngoại lệ, xu hướng.\\[2mm]
\textbf{Giải thích:} truyền đạt kết quả phân tích.\\[2mm]
\textbf{Sản xuất:} tạo hình cho báo cáo, dashboard và ứng dụng.
\end{block}
\end{columns}
\end{frame}

% 7
\begin{frame}[fragile]{Cài đặt và kiểm tra môi trường chỉ cần vài lệnh}
\begin{columns}[T,onlytextwidth]
\column{.52\textwidth}
\begin{lstlisting}[style=py,numbers=none]
# Terminal
python -m pip install matplotlib

# Conda
conda install matplotlib
\end{lstlisting}
\column{.45\textwidth}
\begin{lstlisting}[style=py]
import matplotlib
import matplotlib.pyplot as plt

print(matplotlib.__version__)
\end{lstlisting}
\end{columns}
\vspace{3mm}
\begin{block}{Trong Jupyter Notebook}
Với backend inline mặc định, hình thường xuất hiện ngay dưới cell. \texttt{plt.show()} vẫn nên dùng để mã rõ ràng và chạy ổn định ngoài notebook.
\end{block}
\end{frame}

% 8
\begin{frame}{Một Figure chứa một hoặc nhiều Axes}
\centering
\begin{tikzpicture}[font=\small]
  \draw[very thick,Navy,rounded corners] (0,0) rectangle (12.8,5.2);
  \node[anchor=north west,Blue,font=\bfseries] at (.25,5.0) {Figure: toàn bộ “tờ giấy”};
  \draw[thick,Teal,fill=Paper] (.65,.65) rectangle (6.0,4.3);
  \draw[thick,Orange,fill=white] (6.75,.65) rectangle (12.1,4.3);
  \node[Teal,font=\bfseries] at (3.3,2.55) {Axes 1};
  \node[Orange,font=\bfseries] at (9.4,2.55) {Axes 2};
  \draw[->,Muted] (1.1,1.05) -- (1.1,3.75);
  \draw[->,Muted] (1.1,1.05) -- (5.55,1.05);
  \draw[->,Muted] (7.2,1.05) -- (7.2,3.75);
  \draw[->,Muted] (7.2,1.05) -- (11.65,1.05);
\end{tikzpicture}
\vspace{2mm}

\muted{Axes là vùng vẽ có hệ trục; Axis là từng trục x/y. Không nên dùng lẫn hai thuật ngữ.}
\end{frame}

% 9
\begin{frame}[fragile]{Pyplot nhanh; Object-Oriented rõ ràng và dễ mở rộng}
\begin{columns}[T,onlytextwidth]
\column{.48\textwidth}
\textbf{Pyplot interface}
\begin{lstlisting}[style=py]
plt.plot(x, y)
plt.title("Revenue")
plt.xlabel("Month")
plt.show()
\end{lstlisting}
\good{Phù hợp:} thử nhanh, một biểu đồ đơn giản.

\column{.49\textwidth}
\textbf{Object-Oriented API}
\begin{lstlisting}[style=py]
fig, ax = plt.subplots()
ax.plot(x, y)
ax.set(title="Revenue",
       xlabel="Month")
plt.show()
\end{lstlisting}
\good{Phù hợp:} notebook lớn, subplot, hàm tái sử dụng.
\end{columns}
\vspace{2mm}
\begin{center}\key{Khuyến nghị trong bài này:} ưu tiên Object-Oriented API.\end{center}
\end{frame}

% 10
\begin{frame}[fragile]{Quy trình vẽ ổn định gồm bốn bước}
\begin{columns}[T,onlytextwidth]
\column{.56\textwidth}
\begin{enumerate}\setlength\itemsep{8pt}
  \item Chuẩn bị dữ liệu.
  \item Tạo \texttt{fig, ax}.
  \item Vẽ và tùy biến trên \texttt{ax}.
  \item Hiển thị hoặc lưu \texttt{fig}.
\end{enumerate}
\column{.42\textwidth}
\begin{lstlisting}[style=py]
x = [1, 2, 3, 4]
y = [12, 18, 17, 25]

fig, ax = plt.subplots()
ax.plot(x, y)
ax.set(title="Sales")
fig.tight_layout()
plt.show()
\end{lstlisting}
\end{columns}
\end{frame}

\SectionSlide{2. Năm biểu đồ nền tảng}{Chọn biểu đồ theo câu hỏi phân tích}

% 12
\begin{frame}{Loại câu hỏi quyết định loại biểu đồ}
\small
\begin{tabularx}{\textwidth}{@{}>{\raggedright\arraybackslash}p{3.4cm}>{\raggedright\arraybackslash}p{3.8cm}>{\raggedright\arraybackslash}X@{}}
\toprule
\textbf{Câu hỏi} & \textbf{Biểu đồ} & \textbf{Ví dụ} \\
\midrule
Thay đổi theo thời gian? & Line & Doanh thu theo tháng \\
So sánh các nhóm? & Bar & Doanh số theo khu vực \\
Hai biến liên hệ thế nào? & Scatter & Quảng cáo và doanh số \\
Một biến phân bố ra sao? & Histogram/boxplot & Thời gian giao hàng \\
Cơ cấu tổng thể? & Pie/stacked bar & Tỷ trọng kênh bán \\
Nhiều giá trị trong ma trận? & Heatmap & Tương quan, dữ liệu thiếu \\
\bottomrule
\end{tabularx}
\vspace{1mm}
\footnotesize\muted{Không chọn biểu đồ chỉ vì “trông đẹp”; hãy bắt đầu từ câu hỏi.}
\end{frame}

% 13
\begin{frame}[fragile]{Line chart nhấn mạnh xu hướng và thứ tự thời gian}
\begin{columns}[T,onlytextwidth]
\column{.43\textwidth}
\begin{lstlisting}[style=py]
months = range(1, 13)
revenue = [1.8, 2.0, 2.15,
  2.05, 2.4, 2.55, 2.7,
  2.62, 2.95, 3.1, 3.35, 3.7]

fig, ax = plt.subplots()
ax.plot(months, revenue,
        marker="o")
ax.set(xlabel="Month",
       ylabel="Billion VND")
plt.show()
\end{lstlisting}
\column{.55\textwidth}
\includegraphics[width=\linewidth]{\asset{line_revenue.png}}
\end{columns}
\end{frame}

% 14
\begin{frame}{Khi đọc line chart, hãy tìm bốn loại tín hiệu}
\begin{tightitems}
  \item \key{Mức độ:} giá trị điển hình cao hay thấp?
  \item \key{Xu hướng:} tăng, giảm hay đi ngang?
  \item \key{Mùa vụ:} mẫu hình có lặp lại theo chu kỳ?
  \item \key{Bất thường:} điểm nào lệch khỏi hành vi chung?
\end{tightitems}
\begin{block}{Cẩn thận}
Không nối các nhóm không có thứ tự tự nhiên. Đường nối ngầm khẳng định tính liên tục hoặc trình tự.
\end{block}
\end{frame}

% 15
\begin{frame}[fragile]{Bar chart giúp so sánh độ lớn giữa các nhóm}
\begin{columns}[T,onlytextwidth]
\column{.43\textwidth}
\begin{lstlisting}[style=py]
regions = ["North", "Central",
           "South", "Online"]
sales = [42, 27, 51, 36]

fig, ax = plt.subplots()
bars = ax.bar(regions, sales)
ax.bar_label(bars)
ax.set(ylabel="Billion VND")
plt.show()
\end{lstlisting}
\column{.55\textwidth}
\includegraphics[width=\linewidth]{\asset{bar_region.png}}
\end{columns}
\end{frame}

% 16
\begin{frame}{Bar chart nên bắt đầu trục giá trị từ 0}
\centering\includegraphics[width=.84\textwidth]{\asset{axis_integrity.png}}
\vspace{-2mm}

\begin{itemize}
  \item Chiều dài cột mã hóa độ lớn; cắt trục có thể phóng đại chênh lệch.
  \item Với line chart, trục không nhất thiết bắt đầu từ 0 nếu mục tiêu là xem biến động.
\end{itemize}
\end{frame}

% 17
\begin{frame}[fragile]{Bar ngang hữu ích khi nhãn dài hoặc cần xếp hạng}
\begin{columns}[T,onlytextwidth]
\column{.50\textwidth}
\begin{lstlisting}[style=py]
order = np.argsort(sales)

ax.barh(
    np.array(regions)[order],
    np.array(sales)[order]
)
ax.set_xlabel("Sales")
\end{lstlisting}
\column{.47\textwidth}
\begin{block}{Quy tắc thực hành}
\begin{itemize}
  \item Sắp xếp khi muốn nhấn mạnh thứ hạng.
  \item Giữ thứ tự thời gian hoặc thứ tự nghiệp vụ khi nó có ý nghĩa.
  \item Hạn chế dùng quá nhiều màu nếu màu không mã hóa thông tin.
\end{itemize}
\end{block}
\end{columns}
\end{frame}

% 18
\begin{frame}[fragile]{Scatter plot làm lộ ra quan hệ, cụm và ngoại lệ}
\begin{columns}[T,onlytextwidth]
\column{.43\textwidth}
\begin{lstlisting}[style=py]
fig, ax = plt.subplots()
sc = ax.scatter(
    ad_spend, sales,
    c=sales, cmap="viridis",
    s=55, alpha=.8
)
ax.set(xlabel="Ad spend",
       ylabel="Sales")
fig.colorbar(sc, ax=ax)
\end{lstlisting}
\column{.55\textwidth}
\includegraphics[width=\linewidth]{\asset{scatter_ads_sales.png}}
\end{columns}
\end{frame}

% 19
\begin{frame}{Scatter plot có thể mã hóa thêm biến nhưng dễ quá tải}
\begin{itemize}\setlength\itemsep{7pt}
  \item Vị trí x, y: hai biến số chính.
  \item Màu \texttt{c}: nhóm hoặc biến số thứ ba.
  \item Kích thước \texttt{s}: quy mô, doanh thu hoặc tần suất.
  \item Độ trong suốt \texttt{alpha}: giảm hiện tượng chồng điểm.
\end{itemize}
\begin{block}{Nguyên tắc}
Chỉ thêm một kênh thị giác khi người đọc có thể giải thích nó và legend/colorbar làm rõ ý nghĩa.
\end{block}
\end{frame}

% 20
\begin{frame}[fragile]{Histogram mô tả hình dạng của một phân phối}
\begin{columns}[T,onlytextwidth]
\column{.43\textwidth}
\begin{lstlisting}[style=py]
fig, ax = plt.subplots()
ax.hist(delivery_days,
        bins=18,
        edgecolor="white")
ax.axvline(
    delivery_days.mean(),
    linestyle="--"
)
ax.set(xlabel="Days",
       ylabel="Orders")
\end{lstlisting}
\column{.55\textwidth}
\includegraphics[width=\linewidth]{\asset{hist_delivery.png}}
\end{columns}
\end{frame}

% 21
\begin{frame}{Số bins thay đổi câu chuyện mà histogram kể}
\begin{columns}[T,onlytextwidth]
\column{.47\textwidth}
\begin{block}{Bins quá ít}
Che mất nhiều đỉnh, khoảng trống và ngoại lệ.
\end{block}
\column{.47\textwidth}
\begin{block}{Bins quá nhiều}
Khuếch đại nhiễu ngẫu nhiên và làm hình khó đọc.
\end{block}
\end{columns}
\vspace{4mm}
\begin{itemize}
  \item Thử nhiều giá trị hợp lý thay vì tin một histogram duy nhất.
  \item Có thể dùng \texttt{bins="auto"} làm điểm khởi đầu.
  \item So sánh các nhóm cần dùng chung cạnh bins.
\end{itemize}
\end{frame}

% 22
\begin{frame}[fragile]{Pie chart chỉ phù hợp với cơ cấu đơn giản}
\begin{columns}[T,onlytextwidth]
\column{.43\textwidth}
\begin{lstlisting}[style=py]
sizes = [46, 28, 17, 9]
labels = ["Website", "Store",
          "Marketplace", "Dealer"]

ax.pie(sizes, labels=labels,
       autopct="%1.0f%%",
       startangle=90)
ax.set_title("Revenue share")
\end{lstlisting}
\column{.53\textwidth}
\centering\includegraphics[width=.86\linewidth]{\asset{pie_channel.png}}
\end{columns}
\end{frame}

% 23
\begin{frame}{Pie chart mất hiệu quả khi có nhiều lát hoặc chênh lệch nhỏ}
\begin{columns}[T,onlytextwidth]
\column{.48\textwidth}
\good{Có thể dùng khi}
\begin{itemize}
  \item tổng đúng bằng 100\%;
  \item có ít nhóm;
  \item các tỷ trọng khác biệt rõ.
\end{itemize}
\column{.48\textwidth}
\bad{Nên đổi sang bar khi}
\begin{itemize}
  \item cần so sánh chính xác;
  \item có nhiều nhóm;
  \item có số âm hoặc tổng không rõ.
\end{itemize}
\end{columns}
\vspace{4mm}
\begin{center}\key{Con người so sánh chiều dài tốt hơn góc và diện tích.}\end{center}
\end{frame}

\SectionSlide{3. Tùy biến để biểu đồ tự giải thích}{Nhãn, màu và chú thích phải phục vụ thông điệp}

% 25
\begin{frame}[fragile]{Tiêu đề và nhãn trục là phần của phân tích}
\begin{lstlisting}[style=py]
fig, ax = plt.subplots(figsize=(8, 4.5))
ax.plot(quarters, actual, marker="o", label="Actual")
ax.plot(quarters, target, marker="s", linestyle="--", label="Plan")

ax.set_title("Actual sales exceeded plan in Q4", loc="left")
ax.set_xlabel("Quarter")
ax.set_ylabel("Sales (billion VND)")
ax.legend()
\end{lstlisting}
\vspace{2mm}
\begin{block}{Tiêu đề tốt}
Nêu kết luận chính khi biểu đồ dùng để trình bày; nêu biến và phạm vi khi biểu đồ dùng để khám phá.
\end{block}
\end{frame}

% 26
\begin{frame}[fragile]{Màu, marker và line style phân biệt chuỗi dữ liệu}
\begin{columns}[T,onlytextwidth]
\column{.47\textwidth}
\begin{lstlisting}[style=py]
ax.plot(x, y1,
        color="#1F5AA6",
        marker="o",
        linewidth=2.5)

ax.plot(x, y2,
        color="#F28E2B",
        marker="s",
        linestyle="--")
\end{lstlisting}
\column{.49\textwidth}
\begin{itemize}
  \item Màu: nhấn mạnh hoặc mã hóa nhóm.
  \item Marker: hỗ trợ bản in đen trắng.
  \item Line style: phân biệt thực tế/kế hoạch.
  \item \texttt{alpha}: xử lý chồng lấp.
\end{itemize}
\end{columns}
\end{frame}

% 27
\begin{frame}[fragile]{Legend chỉ hữu ích khi người đọc cần giải mã}
\begin{lstlisting}[style=py]
ax.plot(x, actual, label="Actual")
ax.plot(x, target, label="Plan")
ax.legend(loc="upper left", ncols=2, frameon=False)
\end{lstlisting}
\begin{itemize}
  \item Gán \texttt{label=} ngay khi vẽ mỗi series.
  \item Đặt legend ở vùng trống, tránh che dữ liệu.
  \item Nếu chỉ có một series và tiêu đề/nhãn đã rõ, có thể bỏ legend.
  \item Gắn nhãn trực tiếp vào đường đôi khi tốt hơn legend.
\end{itemize}
\end{frame}

% 28
\begin{frame}[fragile]{Grid và ticks hỗ trợ đọc số nhưng không nên lấn át dữ liệu}
\begin{columns}[T,onlytextwidth]
\column{.48\textwidth}
\begin{lstlisting}[style=py]
ax.grid(axis="y", alpha=.25)
ax.set_xticks(range(4),
              ["Q1", "Q2", "Q3", "Q4"])
ax.tick_params(axis="x",
               rotation=0,
               labelsize=11)
\end{lstlisting}
\column{.48\textwidth}
\begin{block}{Tối giản có chủ đích}
\begin{itemize}
  \item Grid mờ và thường chỉ theo một trục.
  \item Không hiển thị quá nhiều tick labels.
  \item Ghi rõ đơn vị trong nhãn trục.
  \item Dùng formatter cho tiền tệ, phần trăm, ngày tháng.
\end{itemize}
\end{block}
\end{columns}
\end{frame}

% 29
\begin{frame}[fragile]{Annotation biến một điểm dữ liệu thành thông tin hành động}
\begin{columns}[T,onlytextwidth]
\column{.45\textwidth}
\begin{lstlisting}[style=py]
ax.annotate(
    "Above plan",
    xy=(3, actual[3]),
    xytext=(2.1, 124),
    arrowprops={
      "arrowstyle": "->"
    }
)
\end{lstlisting}
\column{.53\textwidth}
\includegraphics[width=\linewidth]{\asset{customized_plot.png}}
\end{columns}
\end{frame}

% 30
\begin{frame}[fragile]{Figure size và layout quyết định khả năng đọc}
\begin{lstlisting}[style=py]
fig, ax = plt.subplots(figsize=(10, 5), constrained_layout=True)
# ... plotting commands ...

# Alternative for an existing figure
fig.tight_layout()
\end{lstlisting}
\begin{itemize}
  \item \texttt{figsize=(width, height)} dùng đơn vị inch.
  \item \texttt{constrained\_layout=True} xử lý khoảng cách tự động ngay khi tạo Figure.
  \item \texttt{tight\_layout()} hữu ích cho bố cục đơn giản.
  \item Luôn kiểm tra hình ở kích thước thực tế của báo cáo/slide.
\end{itemize}
\end{frame}

% 31
\begin{frame}[fragile]{Style sheet tạo tính nhất quán với một lệnh}
\begin{columns}[T,onlytextwidth]
\column{.50\textwidth}
\begin{lstlisting}[style=py]
print(plt.style.available)

plt.style.use("seaborn-v0_8-whitegrid")

with plt.style.context("ggplot"):
    fig, ax = plt.subplots()
    ax.plot(x, y)
\end{lstlisting}
\column{.46\textwidth}
\begin{block}{Cho dự án dài}
Định nghĩa chuẩn màu, font, kích thước và grid bằng \texttt{rcParams} hoặc file \texttt{.mplstyle}; tránh tùy biến lại từng biểu đồ.
\end{block}
\end{columns}
\end{frame}

\SectionSlide{4. Nhiều biểu đồ và đồ thị nâng cao}{Mở rộng từ một Axes sang cấu trúc phân tích}

% 33
\begin{frame}[fragile]{Subplots cho phép so sánh nhiều góc nhìn trong cùng Figure}
\begin{columns}[T,onlytextwidth]
\column{.42\textwidth}
\begin{lstlisting}[style=py]
fig, axes = plt.subplots(
    2, 2,
    figsize=(10, 6),
    constrained_layout=True
)

axes[0, 0].plot(month, revenue)
axes[0, 1].bar(region, sales)
axes[1, 0].hist(delivery)
axes[1, 1].scatter(ads, sales)
\end{lstlisting}
\column{.56\textwidth}
\includegraphics[width=\linewidth]{\asset{subplots_dashboard.png}}
\end{columns}
\end{frame}

% 34
\begin{frame}[fragile]{Chia sẻ trục giúp so sánh công bằng giữa các panels}
\begin{lstlisting}[style=py]
fig, axes = plt.subplots(2, 1, sharex=True, sharey=True)

for ax, region in zip(axes, ["North", "South"]):
    part = df[df["region"] == region]
    ax.plot(part["month"], part["sales"])
    ax.set_title(region)
\end{lstlisting}
\begin{itemize}
  \item Dùng chung scale khi mục tiêu là so sánh mức độ.
  \item Không ép chung scale nếu các đại lượng khác đơn vị hoặc khác bản chất.
  \item Giữ thứ tự panels nhất quán với câu chuyện phân tích.
\end{itemize}
\end{frame}

% 35
\begin{frame}[fragile]{Heatmap giúp quét nhanh ma trận tương quan}
\begin{columns}[T,onlytextwidth]
\column{.41\textwidth}
\begin{lstlisting}[style=py]
corr = df.corr(numeric_only=True)

fig, ax = plt.subplots()
im = ax.imshow(
    corr,
    cmap="RdBu_r",
    vmin=-1, vmax=1
)
fig.colorbar(im, ax=ax)
\end{lstlisting}
\column{.57\textwidth}
\centering\includegraphics[width=.80\linewidth]{\asset{correlation_heatmap.png}}
\end{columns}
\end{frame}

% 36
\begin{frame}{Heatmap hỗ trợ phát hiện tương quan và cấu trúc thiếu dữ liệu}
\begin{columns}[T,onlytextwidth]
\column{.55\textwidth}
\includegraphics[width=\linewidth]{\asset{missing_heatmap.png}}
\column{.42\textwidth}
\begin{itemize}
  \item Cột nhiều ô thiếu: cân nhắc bỏ biến hoặc tìm nguồn thay thế.
  \item Hàng thiếu theo cụm: kiểm tra quy trình thu thập.
  \item Mẫu thiếu có cấu trúc: không nên mặc định là ngẫu nhiên.
\end{itemize}
\end{columns}
\vspace{2mm}
\small\muted{Heatmap giúp quét nhanh; không thay thế biểu đồ cần đọc giá trị chính xác.}
\end{frame}

% 37
\begin{frame}[fragile]{Contour thường dễ đọc hơn surface 3D}
\begin{columns}[T,onlytextwidth]
\column{.42\textwidth}
\begin{lstlisting}[style=py]
X, Y = np.meshgrid(x, y)
Z = f(X, Y)

fig = plt.figure()
ax3d = fig.add_subplot(
    121, projection="3d")
ax3d.plot_surface(X, Y, Z,
                  cmap="viridis")

ax2d = fig.add_subplot(122)
ax2d.contourf(X, Y, Z)
\end{lstlisting}
\column{.56\textwidth}
\includegraphics[width=\linewidth]{\asset{surface_contour.png}}
\end{columns}
\end{frame}

% 38
\begin{frame}{3D chỉ nên dùng khi chiều thứ ba thật sự quan trọng}
\begin{columns}[T,onlytextwidth]
\column{.48\textwidth}
\good{Dùng 3D cho}
\begin{itemize}
  \item bề mặt toán học;
  \item địa hình;
  \item mô hình không gian;
  \item tương tác xoay góc nhìn.
\end{itemize}
\column{.48\textwidth}
\bad{Tránh 3D cho}
\begin{itemize}
  \item bar/pie trang trí;
  \item so sánh chính xác;
  \item hình tĩnh bị che khuất;
  \item “làm biểu đồ ấn tượng”.
\end{itemize}
\end{columns}
\end{frame}

% 39
\begin{frame}[fragile]{Animation mô tả trạng thái thay đổi theo thời gian}
\begin{lstlisting}[style=py]
from matplotlib.animation import FuncAnimation

fig, ax = plt.subplots()
line, = ax.plot([], [])

def update(frame):
    line.set_data(x[:frame], y[:frame])
    return line,

ani = FuncAnimation(fig, update, frames=len(x), interval=80)
ani.save("trend.gif", writer="pillow")
\end{lstlisting}
\muted{Animation phù hợp khi chuyển động là thông tin; không nên thay một biểu đồ tĩnh rõ ràng bằng hiệu ứng.}
\end{frame}

% 40
\begin{frame}{Widgets biến Figure thành công cụ khám phá nhỏ}
\begin{itemize}\setlength\itemsep{7pt}
  \item \texttt{matplotlib.widgets.Slider}: thay đổi ngưỡng hoặc tham số.
  \item \texttt{Button}: chạy lại phép tính hoặc đổi chế độ.
  \item \texttt{CheckButtons}: bật/tắt series.
  \item \texttt{RangeSlider}: lọc khoảng giá trị.
\end{itemize}
\begin{block}{Giới hạn}
Matplotlib có tương tác, nhưng dashboard web nhiều người dùng thường phù hợp hơn với Plotly, Bokeh, Streamlit hoặc Dash.
\end{block}
\end{frame}

\SectionSlide{5. Pandas, Seaborn và xuất bản}{Từ phân tích nhanh đến hình dùng trong báo cáo}

% 42
\begin{frame}[fragile]{Pandas cung cấp API ngắn gọn nhưng vẫn trả về Axes}
\begin{columns}[T,onlytextwidth]
\column{.42\textwidth}
\begin{lstlisting}[style=py]
ax = df.plot(
    y=["revenue", "cost"],
    marker="o",
    figsize=(8, 4)
)

ax.set_title("Revenue and cost")
ax.set_ylabel("Billion VND")
ax.legend(frameon=False)
\end{lstlisting}
\column{.56\textwidth}
\includegraphics[width=\linewidth]{\asset{pandas_plot.png}}
\end{columns}
\end{frame}

% 43
\begin{frame}{Matplotlib và Seaborn bổ sung cho nhau}
\small
\begin{tabularx}{\textwidth}{@{}p{3.5cm}XX@{}}
\toprule
 & \textbf{Matplotlib} & \textbf{Seaborn} \\
\midrule
Mức trừu tượng & Thấp hơn, kiểm soát chi tiết & Cao hơn, thiên về thống kê \\
Dữ liệu & Mảng, chuỗi, DataFrame & DataFrame dạng tidy \\
Thế mạnh & Tùy biến, layout, xuất bản & Phân phối, nhóm, ước lượng \\
Kết quả & Figure/Axes & Vẫn dựa trên Matplotlib \\
\bottomrule
\end{tabularx}
\vspace{4mm}
\begin{center}\key{Cách làm phổ biến:} Seaborn vẽ nhanh, Matplotlib hoàn thiện tiêu đề, trục và bố cục.\end{center}
\end{frame}

% 44
\begin{frame}[fragile]{Kết hợp Seaborn và Matplotlib không làm mất quyền kiểm soát}
\begin{lstlisting}[style=py]
import seaborn as sns
import matplotlib.pyplot as plt

fig, ax = plt.subplots(figsize=(8, 4))
sns.boxplot(data=df, x="segment", y="order_value", ax=ax)
ax.set(title="Order value by customer segment",
       xlabel="Customer segment", ylabel="Order value (VND)")
fig.tight_layout()
\end{lstlisting}
\begin{block}{Ghi nhớ}
Truyền rõ \texttt{ax=ax} giúp kiểm soát Figure và tích hợp biểu đồ vào subplot.
\end{block}
\end{frame}

% 45
\begin{frame}[fragile]{savefig xuất hình cho slide, web và bài báo}
\begin{lstlisting}[style=py]
fig.savefig("revenue.png", dpi=300, bbox_inches="tight")
fig.savefig("revenue.pdf", bbox_inches="tight")
fig.savefig("revenue.svg", bbox_inches="tight")
\end{lstlisting}
\begin{columns}[T,onlytextwidth]
\column{.31\textwidth}\textbf{PNG/JPG}\\Ảnh raster; phù hợp web, slide.
\column{.31\textwidth}\textbf{PDF/SVG}\\Vector; nét khi phóng to, phù hợp bài báo.
\column{.31\textwidth}\textbf{DPI}\\Quan trọng với raster; 150--300 thường đủ.
\end{columns}
\vspace{4mm}
\muted{Lưu trước \texttt{plt.show()} để tránh khác biệt giữa các backend.}
\end{frame}

% 46
\begin{frame}[fragile]{Một PDF có thể chứa nhiều Figure}
\begin{lstlisting}[style=py]
from matplotlib.backends.backend_pdf import PdfPages

with PdfPages("report_figures.pdf") as pdf:
    for region in regions:
        fig, ax = plt.subplots()
        plot_region(ax, df, region)
        pdf.savefig(fig, bbox_inches="tight")
        plt.close(fig)
\end{lstlisting}
\begin{block}{Thực hành tốt}
Đóng Figure trong vòng lặp bằng \texttt{plt.close(fig)} để tránh tăng bộ nhớ khi tạo hàng trăm hình.
\end{block}
\end{frame}

% 47
\begin{frame}{Các toolkit mở rộng Matplotlib theo loại dữ liệu}
\begin{tightitems}
  \item \key{mplot3d:} 3D tích hợp trong Matplotlib.
  \item \key{Seaborn:} biểu đồ thống kê cấp cao.
  \item \key{GeoPandas:} bản đồ từ GeoDataFrame, dùng Matplotlib làm backend.
  \item \key{Cartopy:} phép chiếu và dữ liệu địa lý chuyên sâu.
  \item \key{mplfinance:} biểu đồ tài chính.
\end{tightitems}
\begin{block}{Nguyên tắc chọn công cụ}
Dùng lớp trừu tượng cao để làm nhanh; quay về Figure/Axes khi cần kiểm soát bố cục và xuất bản.
\end{block}
\end{frame}

\SectionSlide{6. Matplotlib trong quy trình Data Science}{Biểu đồ là công cụ kiểm tra giả định và truyền đạt kết quả}

% 49
\begin{frame}{Trực quan hóa xuất hiện xuyên suốt vòng đời phân tích}
\centering
\begin{tikzpicture}[node distance=1.55cm,>=stealth,font=\small]
\tikzset{processbox/.style={draw=Blue,very thick,rounded corners,fill=Paper,minimum width=2.25cm,minimum height=1.05cm,align=center}}
\node[processbox] (q) {Câu hỏi\\nghiệp vụ};
\node[processbox,right of=q,xshift=1.15cm] (clean) {Làm sạch\\dữ liệu};
\node[processbox,right of=clean,xshift=1.15cm] (eda) {Khám phá\\EDA};
\node[processbox,below of=eda,yshift=-.45cm] (model) {Mô hình\\hóa};
\node[processbox,left of=model,xshift=-1.15cm] (eval) {Đánh giá\\mô hình};
\node[processbox,left of=eval,xshift=-1.15cm] (comm) {Truyền đạt\\quyết định};
\draw[->,thick] (q)--(clean); \draw[->,thick] (clean)--(eda);
\draw[->,thick] (eda)--(model); \draw[->,thick] (model)--(eval);
\draw[->,thick] (eval)--(comm); \draw[->,thick] (comm)--(q);
\end{tikzpicture}
\vspace{5mm}

\muted{EDA không phải “bước trang trí”; nó hỗ trợ phát hiện lỗi, chọn biến, biến đổi dữ liệu và hình thành giả thuyết.}
\end{frame}

% 50
\begin{frame}{Mục tiêu học máy định hướng cách trực quan hóa}
\small
\begin{tabularx}{\textwidth}{@{}p{3.1cm}X@{}}
\toprule
\textbf{Mục tiêu} & \textbf{Cách trực quan hóa hữu ích} \\
\midrule
Dự báo số & Đặt target ở trục y; scatter với predictor; histogram/boxplot để xem biến đổi \\
Phân loại & Boxplot predictor theo class; scatter tô màu theo class; confusion matrix \\
Chuỗi thời gian & Line chart ở nhiều mức tổng hợp; zoom vào giai đoạn bất thường \\
Không giám sát & Scatter matrix; heatmap tương quan; biểu đồ profile của cụm \\
Chất lượng dữ liệu & Histogram, boxplot, missing-value heatmap, kiểm tra giá trị bất hợp lệ \\
\bottomrule
\end{tabularx}
\end{frame}

% 51
\begin{frame}{Sáu câu hỏi kiểm tra trước khi công bố biểu đồ}
\begin{enumerate}\setlength\itemsep{6pt}
  \item Biểu đồ trả lời câu hỏi nào?
  \item Kiểu biểu đồ có phù hợp kiểu dữ liệu?
  \item Trục, đơn vị và phạm vi có trung thực?
  \item Màu có mang ý nghĩa hay chỉ trang trí?
  \item Người đọc có thấy thông điệp trong 5--10 giây?
  \item Mã và dữ liệu có đủ để tái lập hình?
\end{enumerate}
\end{frame}

% 52
\begin{frame}{Các lỗi phổ biến làm giảm độ tin cậy}
\begin{columns}[T,onlytextwidth]
\column{.48\textwidth}
\bad{Lỗi thiết kế}
\begin{itemize}
  \item trục bị cắt gây phóng đại;
  \item quá nhiều màu/legend;
  \item font và nhãn quá nhỏ;
  \item 3D che khuất dữ liệu;
  \item pie chart có quá nhiều lát.
\end{itemize}
\column{.48\textwidth}
\bad{Lỗi kỹ thuật}
\begin{itemize}
  \item dùng stateful API trong code phức tạp;
  \item quên đóng Figure trong vòng lặp;
  \item lưu ảnh DPI thấp;
  \item layout cắt nhãn;
  \item không cố định seed cho dữ liệu mô phỏng.
\end{itemize}
\end{columns}
\end{frame}

% 53
\begin{frame}[fragile]{Đóng gói biểu đồ thành hàm giúp tái sử dụng và kiểm thử}
\begin{lstlisting}[style=py]
def plot_monthly_sales(ax, data, title):
    ax.plot(data["month"], data["sales"], marker="o")
    ax.set(title=title, xlabel="Month", ylabel="Sales")
    ax.grid(axis="y", alpha=.25)
    return ax

fig, ax = plt.subplots(figsize=(8, 4))
plot_monthly_sales(ax, df, "Monthly sales")
fig.savefig("monthly_sales.svg", bbox_inches="tight")
\end{lstlisting}
\muted{Hàm nhận \texttt{ax} dễ ghép vào subplot và không phụ thuộc trạng thái toàn cục.}
\end{frame}

\SectionSlide{7. Củng cố và thực hành}{Kiểm tra hiểu biết rồi áp dụng vào dữ liệu kinh doanh}

% 55
\begin{frame}{Quiz 1: Chọn biểu đồ phù hợp}
Bạn cần so sánh phân phối giá trị đơn hàng giữa bốn phân khúc khách hàng. Biểu đồ nào phù hợp nhất?
\begin{enumerate}[A.]\setlength\itemsep{8pt}
  \item Line chart
  \item Side-by-side boxplot
  \item Pie chart
  \item 3D surface
\end{enumerate}
\pause
\begin{block}{Đáp án: B}
Boxplot cho thấy median, quartiles, độ phân tán và ngoại lệ của từng phân khúc.
\end{block}
\end{frame}

% 56
\begin{frame}{Quiz 2: Figure và Axes}
Trong lệnh \texttt{fig, ax = plt.subplots()}, phát biểu nào đúng?
\begin{enumerate}[A.]\setlength\itemsep{7pt}
  \item \texttt{fig} là trục x; \texttt{ax} là trục y.
  \item \texttt{fig} là toàn bộ canvas; \texttt{ax} là vùng vẽ.
  \item \texttt{fig} chỉ dùng để lưu PNG.
  \item \texttt{ax} chỉ dùng khi có nhiều subplot.
\end{enumerate}
\pause
\begin{block}{Đáp án: B}
Một Figure có thể chứa một hoặc nhiều Axes; Axes vẫn được dùng ngay cả khi chỉ có một biểu đồ.
\end{block}
\end{frame}

% 57
\begin{frame}{Quiz 3: Phát hiện vấn đề trình bày}
Một bar chart so sánh doanh số 98 và 100, nhưng trục y bắt đầu từ 97. Vấn đề chính là gì?
\pause
\begin{block}{Đáp án}
Chiều dài cột bị diễn giải từ mốc 0, nên trục cắt làm chênh lệch 2\% trông rất lớn. Nếu cần soi biến động nhỏ, dùng dot plot hoặc line chart kèm giá trị và giải thích phạm vi trục.
\end{block}
\end{frame}

% 58
\begin{frame}[fragile]{Bài tập mẫu: phân tích doanh thu theo tháng}
\textbf{Yêu cầu:} vẽ line chart, thêm marker, đường trung bình và chú thích tháng cao nhất.
\begin{lstlisting}[style=py]
months = np.arange(1, 13)
revenue = np.array([18, 20, 22, 21, 24, 26, 27, 26, 30, 31, 34, 37])

fig, ax = plt.subplots(figsize=(9, 4.5))
ax.plot(months, revenue, marker="o")
ax.axhline(revenue.mean(), linestyle="--", color="gray", label="Mean")
i = revenue.argmax()
ax.annotate("Peak", xy=(months[i], revenue[i]), xytext=(9, 38),
            arrowprops={"arrowstyle": "->"})
ax.set(xlabel="Month", ylabel="Revenue", title="Monthly revenue")
ax.legend(); fig.tight_layout()
\end{lstlisting}
\end{frame}

% 59
\begin{frame}{Bài thực hành 1: Dashboard bán hàng 2×2}
Từ DataFrame có các cột \texttt{date, region, channel, revenue, cost, orders}:
\begin{enumerate}\setlength\itemsep{5pt}
  \item Tính doanh thu theo tháng và vẽ line chart.
  \item Tính doanh thu theo vùng và vẽ bar chart đã sắp xếp.
  \item Vẽ histogram giá trị đơn hàng trung bình.
  \item Vẽ scatter giữa chi phí và doanh thu, tô màu theo kênh.
  \item Ghép bốn biểu đồ vào một Figure 2×2, thêm tiêu đề chung.
  \item Lưu kết quả thành PNG 300 DPI và PDF.
\end{enumerate}
\textbf{Gợi ý:} \texttt{groupby}, \texttt{resample}, \texttt{plt.subplots}, \texttt{fig.suptitle}.
\end{frame}

% 60
\begin{frame}{Bài thực hành 2: Chẩn đoán dữ liệu giao hàng}
Với dữ liệu gồm \texttt{distance, weight, promised\_days, actual\_days, carrier}:
\begin{enumerate}\setlength\itemsep{5pt}
  \item Tạo biến \texttt{delay = actual\_days - promised\_days}.
  \item Dùng histogram để xem phân phối delay.
  \item Dùng boxplot để so sánh delay giữa các hãng vận chuyển.
  \item Dùng scatter để khảo sát distance và actual\_days.
  \item Chú thích ba đơn hàng trễ nhất.
  \item Viết 3 nhận xét nghiệp vụ dựa trên biểu đồ.
\end{enumerate}
\end{frame}

% 61
\begin{frame}{Bài tập tự làm: sửa một biểu đồ gây hiểu sai}
Tìm hoặc tự tạo một biểu đồ có ít nhất ba vấn đề:
\begin{itemize}
  \item lựa chọn loại biểu đồ chưa phù hợp;
  \item trục hoặc đơn vị không rõ;
  \item màu/legend dư thừa;
  \item thiếu tiêu đề mang thông điệp;
  \item nhãn bị chồng hoặc quá nhỏ.
\end{itemize}
Nộp: (1) hình ban đầu, (2) hình đã sửa, (3) đoạn giải thích 150--200 từ, (4) mã Python có thể chạy lại.
\end{frame}

% 62
\begin{frame}{Tóm tắt: từ câu hỏi đến quyết định}
\footnotesize
\begin{tightitems}
  \item \key{Figure--Axes} là mô hình tư duy cốt lõi; ưu tiên Object-Oriented API.
  \item Line, bar, scatter và histogram giải quyết phần lớn nhu cầu EDA cơ bản.
  \item Tùy biến chỉ có giá trị khi làm rõ dữ liệu hoặc thông điệp.
  \item Subplots, heatmap, 3D và animation phục vụ những cấu trúc câu hỏi cụ thể.
  \item Pandas/Seaborn giúp làm nhanh; Matplotlib giữ quyền kiểm soát cuối cùng.
  \item Xuất bản cần kiểm tra trục, đơn vị, layout, định dạng và khả năng tái lập.
\end{tightitems}
\end{frame}

% 63
\begin{frame}{Tài liệu tham khảo và học tiếp}
\small
\begin{thebibliography}{9}
\bibitem{mpl} Matplotlib Development Team. \emph{Matplotlib Documentation}.\\
\url{https://matplotlib.org/stable/}

\bibitem{gfg} GeeksforGeeks. \emph{Matplotlib Tutorial}, cập nhật 24/02/2026.\\
\url{https://www.geeksforgeeks.org/python-matplotlib-an-overview/}

\bibitem{shmueli} G. Shmueli, P. C. Bruce, A. V. Deokar, N. R. Patel (2023).\\
\emph{Machine Learning for Business Analytics}, Chương 3: Data Visualization. Wiley.

\bibitem{few} S. Few (2012). \emph{Show Me the Numbers}, 2nd ed. Analytics Press.
\end{thebibliography}
\end{frame}



\end{document}
